% 表紙・凡例・目次（記入例版）
\thispagestyle{cover}
\vspace*{12mm}
{\sffamily\color{pone}\large 経営情報 2026　AI利用ガイド}\par
\vspace{3mm}
{\sffamily\bfseries\fontsize{30pt}{38pt}\selectfont 復習教材}\par
\vspace{2mm}
{\sffamily\fontsize{20pt}{26pt}\selectfont 問いを立て、調べ、読み、磨く　\textcolor{ans}{記入例}}\par
\vspace{3mm}
{\color{inktwo}\rule{\linewidth}{0.6pt}}\par
\vspace{6mm}

\begin{tcolorbox}[enhanced,colback=papertwo,colframe=papertwo,boxrule=0pt,arc=1.5mm,
    left=4mm,right=4mm,top=3mm,bottom=3mm,fontupper=\small]
{\sffamily\bfseries この記入例の見方}\par\vspace{1mm}
\begin{itemize}[label=・]
  \item ワークシートに、\ans{赤字}で書き込みを入れています。まず自分で取り組み、その後に見比べてください。
  \item 答えが一つに定まる設問（良い問いの条件など）の書き込みは\textbf{模範解答}です。
  \item 人によって答えが変わる設問（テーマ、問い、文献、要約など）の書き込みは\textbf{記入例}です。下の架空の学生が書いたものとして、5つの単元を一貫させています。正解ではなく、書く量と具体性の目安です。
  \item 各単元の末尾の「\textcolor{ans}{解説・ポイント}」に、取り組むときの要点とつまずきやすい点をまとめています。
\end{itemize}
\end{tcolorbox}

\vspace{4mm}
\begin{tcolorbox}[enhanced,colback=white,colframe=rulegray,boxrule=0.5pt,arc=1.5mm,
    left=4mm,right=4mm,top=3mm,bottom=3mm,fontupper=\small]
{\sffamily\bfseries 記入例の前提（架空）}\par\vspace{1mm}
\begin{itemize}[label=・]
  \item \textbf{学生}：経営情報の受講生。アルバイト先の町工場で、社長が「生成AIを使ってみたいが、何に使えばよいかわからない」と話していたことから、中小企業と生成AIの関係に関心を持った。
  \item \textbf{問い}：復習2で「愛知県の中小企業では、なぜ生成AIの業務利用が進まないのか」を仮の問いとし、復習5で「生成AIを試験導入した愛知県の中小製造業は、どの業務で使い、本格導入の障壁を何と捉えているのか」に修正して確定する。
  \item \textbf{文献}：記入例に挙げた白書・統計・ガイドラインは実在する継続的な資料です。版や年次、URLは各自で確認してください。復習4で精読する論文は、記入例のための\textbf{架空の論文}です。
  \item \textbf{AIの出力}：記入例に書いたAIの出力は一例です。実際の出力は、使うたびに変わります。
\end{itemize}
\end{tcolorbox}

\vspace{6mm}
{\sffamily\bfseries 目次}\par\vspace{1mm}
{\small\makeatletter\@starttoc{toc}\makeatother}
